\section{Methodology}
\label{sec:method}

\subsection{State Representation}
\label{sec:state}

\textbf{Position.}
The world-frame position $\vp_w(t) \in \mathbb{R}^3$ is a quintic
(order-6) uniform B-spline with knot spacing $\Delta t_p = 40$\,ms.
Velocity $\dot{\vp}_w$ and acceleration $\ddot{\vp}_w$ are
obtained as analytical derivatives of the spline.

\textbf{Orientation.}
The body-to-world rotation $\vR(t)=\vR_{wb}(t)$ is a cubic
cumulative $\SO$ B-spline~\cite{sommer2020efficient,usenko2020basalt}:
\begin{equation}
  \vR(t) = \vR_{k-3}
           \prod_{j=k-2}^{k} \Exp\!\bigl(\tilde{B}_j(t)\,
           \Log(\vR_{j-1}^{-1}\vR_{j})\bigr),
  \label{eq:ori_spline}
\end{equation}
where $\vR_{k-3},\ldots,\vR_k$ are the four active orientation control points
and $\tilde{B}_j(t)$ are the cumulative basis functions defined in
\autoref{sec:prelim_so3}.  The absolute orientation control points are optimized.
The knot spacing is $\Delta t_o = 16$\,ms, and the body-frame angular
velocity $\vom_b(t)$ follows by analytic differentiation.
Both spline spacings are fixed across flights and platforms
(\autoref{sec:negative}).

\textbf{Biases and extrinsics.}
The constant six-component \ac{imu} bias $[\vb_a,\vb_g]$ is estimated
jointly with the trajectory.  A Gaussian prior anchors it to the
stationary initialization (\autoref{sec:init}), with standard
deviations $0.216$\,m/s$^2$ and $0.0036$\,rad/s for the accelerometer
and gyroscope, respectively.  These widths are measured offline from stationary ground segments,
the accelerometer's against the specific force $-\vR_{bw}\vg_w$
predicted from the \ac{mocap} attitude, the gyroscope's from the
bias shift between motors off and motors running.  The prior weights are inverse variances, and the prior touches no leaving state, so it is re-applied in every window and never absorbed into the marginalization prior (\autoref{sec:sw}).
All radar extrinsics are fixed at the platform's measured values.
On our platform, freeing the pitch left the estimate close to its
initial value, so we use the measured $27.5\degree$ mount angle
(\autoref{sec:overview}).

\subsection{Measurement Models}
\label{sec:models}

\textbf{Radar Doppler with lever arm.}
For a static target at sensor-frame bearing $\vu_s$, the predicted
Doppler is
\begin{equation}
  v_D = -\vu_b^{\top}\!\left(
    \vR_{bw}\,\vv_w(t) + \vom_b(t) \times \mathbf{t}_{bs}
  \right),
  \label{eq:doppler}
\end{equation}
where $\vv_w=\dot{\vp}_w$, $\vu_b=\vR_{bs}\vu_s$,
$\vR_{bs}$ rotates sensor coordinates into the body frame, and
$\vR_{bw}=\vR_{wb}^{\top}$.  The lever arm $\mathbf{t}_{bs}$ runs
from the body \ac{com} to the radar antenna, expressed in body coordinates.
The sign follows the sensor convention: positive Doppler denotes a
receding target.  Each return contributes a residual
$r_{f,j}=v_{D,\mathrm{meas}}-v_{D,\mathrm{pred}}$ to its frame's
jointly weighted residual block (\autoref{sec:robust}).

\textbf{Accelerometer.}
\begin{equation}
  \mathbf{r}_a = \vz_a - \vR_{bw}\!\left(\ddot{\vp}_w - \vg_w\right) - \vb_a,
  \label{eq:accel}
\end{equation}
where $\vz_a$ is the measured specific force and $\vg_w$ is gravity
in world coordinates.

\textbf{Gyroscope.}
\begin{equation}
  \mathbf{r}_g = \vz_g - \vom_b(t) - \vb_g,
  \label{eq:gyro}
\end{equation}
Here $\vz_g$ is the measured body-frame angular rate.  The weights
for both \ac{imu} residuals are specified in \autoref{sec:robust}.

\textbf{Heading prior.}
Absolute yaw is unobservable from Doppler and \ac{imu} measurements:
a common rotation of position, velocity and attitude about gravity
leaves these measurements unchanged.  We supply heading using
\ac{mocap} yaw at approximately 100\,Hz as an idealized magnetometer
surrogate.  With $\mathbf{h} = \vR\,\mathbf{e}_x$ the body forward axis in world
coordinates, the residual is
\begin{equation}
  r_\psi = h_x \sin\psi_{\mathrm{ref}} - h_y \cos\psi_{\mathrm{ref}},
  \label{eq:heading}
\end{equation}
i.e.\ $\sin(\psi-\psi_{\mathrm{ref}})$ scaled by the horizontal
projection of the forward axis, with weight $\lambda_\psi=0.6$.  It
needs no angle wrapping, vanishes when the nose points vertically (the
heading is momentarily undefined through a flip), and has a spurious
zero at the opposite heading, which the initialization avoids.
\autoref{sec:yaw} discusses this weight and the limitations of the surrogate.

\subsection{Radar Front End and Measurement Weighting}
\label{sec:robust}

The radar front end unwraps Doppler, rejects consensus outliers and
assigns weights before optimization.  The measurement characterization
in \autoref{sec:char} motivates these steps.

\textbf{Alias unwrapping and the alias weight.}
At frame load, each return's Doppler alias is selected using the
radial velocity predicted from \ac{imu} integration
(\autoref{sec:init}).  Correct selection requires the predicted radial
velocity to be within $v_{\max}$ of the true value.  All subsequent
stages use this unwrapped Doppler.
Unwrapping does not correct the bearing errors of aliased returns.
These returns retain their reported bearings but receive the reduced
weight $w_{\mathrm{al}}=0.019$, the measured ratio of clean to aliased
noise-floor variances (\autoref{sec:char_alias}).

\textbf{Consensus prefilter.}
Following \mbox{reve}~\cite{doer2020ekf}, three-point \ac{ransac}
estimates ego-velocity for each frame.  Only returns within
$0.15$\,m/s of the consensus prediction reach the solver
(\autoref{sec:char_outliers}).  Frames with fewer than five returns
bypass this filter.  Filtering runs once at frame load.

\textbf{Per-return weights in the frame covariance.}
For each retained, non-aliased return, the relative inverse-variance
weight from \eqref{eq:law_fit} is, in the simplified angular model,
\begin{equation}
  w_{f,j} = \frac{1}{1 + (c\,q_{f,j}/\sigma_0)^2}, \qquad
  q_{f,j} = \frac{\lVert\vv_f\rVert\sin\theta_j}{\cos\phi_j},
  \label{eq:weight}
\end{equation}
where $\theta_j$ is the angle between velocity and ray, and $\phi_j$
is the angle off boresight.  The constants are
$c=\delta\phi_0/\sqrt2$, $\delta\phi_0=2.7\degree$ (converted to
radians), and $\sigma_0=0.083$\,m/s (\autoref{sec:char_rate}).
The solver uses the direction-cosine form of the same law: with
$\mathbf{s}$ the unit bearing and $\vv_s$ the frame velocity, both in
sensor coordinates, and $c_\phi = \max(s_x, 0.15)$,
\begin{equation}
\begin{aligned}
  X_j &= \sqrt{(v_{s,y} - s_y v_{s,x}/c_\phi)^2 + (v_{s,z} - s_z v_{s,x}/c_\phi)^2},\\
  \sigma_j^2 &= \sigma_0^2 + (c\,X_j)^2,
\end{aligned}
\label{eq:dircos}
\end{equation}
which needs no clamp towards $\phi = 90\degree$ beyond the floor on
$c_\phi$; $w_{f,j} = \sigma_0^2/\sigma_j^2$.
The velocity $\vv_f$ comes from the frame's \ac{wls} consensus fit.
Weights are fixed when the problem is built, independently of later
trajectory updates.  They still depend on the \ac{imu} prediction
through the earlier alias decisions.

The residuals of frame $f$ are stacked in $\mathbf{r}_f$.  Their
covariance model is
\begin{equation}
  \Sigma_f = \sigma_0^2\!\left[
    \mathrm{diag}(w_{f,j}^{-1})
    + \sigma_c^2\mathbf{1}\mathbf{1}^{\top}\right],
  \label{eq:frame_cov}
\end{equation}
using $w_{f,j}=w_{\mathrm{al}}$ for aliased returns and
$\sigma_c=1.27$ for the shared component (\autoref{sec:char_corr}).
The rank-one term accounts for a common additive error within a frame,
reducing the information attributed to repeated observations of that
component.  The frame is weighted jointly, as specified below.

\textbf{Accelerometer gate.}
The accelerometer weight decreases with body-rate magnitude to limit
the influence of vibration and other unmodeled high-dynamic errors.
The gate in \eqref{eq:comp_a} uses an empirical scale
$\omega_a=8$\,rad/s, chosen from the velocity--orientation trade-off.
Translational acceleration itself is already included in
\eqref{eq:accel}.  The gyroscope weight remains constant.

\textbf{Static stream weights.}
The per-sample stream weights follow the correlation-adjusted error
budgets of \eqref{eq:lambda_star}.  Taking their geometric means over
the two racing flights gives $\lambda_g=3.13$, $\lambda_a=0.0039$
and $\lambda_r=2.64$.  These weights are fixed across flights and
platforms.

\textbf{The complete weighting.}
The effective information weights are
\begin{align}
  \Lambda^{\mathrm{gyro}}  &= \lambda_g, \label{eq:comp_g}\\
  \Lambda^{\mathrm{accel}} &= \frac{\lambda_a}{1 + (\lVert\vom_b\rVert/\omega_a)^2}, \label{eq:comp_a}\\
  \Lambda^{\mathrm{radar}}_f &= \lambda_r\,\sigma_0^2\Sigma_f^{-1}. \label{eq:comp_r}
\end{align}
Thus the radar contribution is
$\sum_f\mathbf{r}_f^{\top}\Lambda^{\mathrm{radar}}_f\mathbf{r}_f$.
The factor $\sigma_0^2$ normalizes the covariance because $\lambda_r$
already includes the stream's inverse-variance scale.  Without
within-frame correlation, this reduces to
$\lambda_r\sum_{f,j}w_{f,j}r_{f,j}^2$.
All retained measurement residuals enter quadratically, without a
robust kernel.  The constants in these weights are measured except
for the empirical accelerometer gate.

\subsection{Regularization}
\label{sec:reg}

\textbf{Position: minimum snap.}
$E_{\text{snap}} = \lambda_s \sum_i \|\vp^{(4)}(t_i + \Delta t_p/2)\|^2$,
one term per knot interval (the midpoint sum of the snap integral,
without the $\Delta t_p$ factor), with $\lambda_s = 2{\times}10^{-5}$.

\textbf{Orientation: angular-increment smoothness.}
We penalize changes between consecutive rotation increments:
\begin{equation}
  \mathbf{r}_\alpha = \Log(\vR_{i-1}^{-1}\vR_i)
                    - \Log(\vR_i^{-1}\vR_{i+1}),
  \label{eq:min_alpha}
\end{equation}
The cost is $\lambda_\alpha\sum_i\lVert\mathbf{r}_{\alpha,i}\rVert^2$
with $\lambda_\alpha=0.1$.  This discrete smoothness penalty is zero
for equal consecutive increments.  It approximates angular-acceleration
regularization at the fixed knot spacing but is not a time-normalized
angular-acceleration residual.  Both regularization weights are design
choices (\autoref{sec:fixed_constants}).

\subsection{Sensor-Only Initialization (P1--P3)}
\label{sec:init}

\textbf{P1: Orientation.}
An automatically detected stationary segment provides the mean
gyroscope bias and a gravity-aligned scalar accelerometer correction.
This does not identify the full accelerometer bias independently of
tilt.  Initial roll and pitch follow from the corrected gravity
direction, and yaw is set to zero until heading aiding is applied.
Orientation is initialized by integrating debiased gyroscope readings:
$\vR(t+\Delta t) = \vR(t)\,\Exp\!\bigl((\vz_g - \vb_g)\Delta t\bigr)$.

\textbf{P2: Position.}
The front end of \autoref{sec:robust} produces a per-frame \ac{wls}
ego-velocity estimate from the unwrapped, prefiltered returns.
Rotated into the world frame with the initialized attitude (the lever
arm is neglected at this stage), it is integrated with forward Euler
into the initial position trajectory.  The alias predictor of
\autoref{sec:robust} is the \ac{imu}-integrated velocity from the same
attitude, and a frame whose consensus fails keeps all its returns.

\textbf{P3: Boundary priors.}
Position, velocity and orientation priors anchor the first window at
$t=0$ to the initialized state, with hand-set weight
$\lambda_{\text{bnd}}=1000$.  They are not re-applied in later windows,
which carry previous information through the marginalization prior
(\autoref{sec:sw}).

\subsection{Sliding Window with Marginalization}
\label{sec:sw}

The fixed-lag smoother advances a 3\,s window in 0.3\,s strides.
After each solve, position and orientation control points leaving the
window are eliminated from the linearized problem by a Schur
complement.  The retained variables coupled to these states, including
the shared bias, form the boundary of the next window.  The resulting
information matrix is stored as a Gaussian prior, with the
approximations described below.

\textbf{Which factors enter the marginal.}
Factors touching a leaving variable, including the previous prior
when applicable, enter the Schur complement and are then removed from
the active problem.  Factors involving only retained variables remain
active and must not also be absorbed into the prior.  In particular,
sharing the retained \ac{imu} bias does not justify including every
\ac{imu} factor.  The boundary is the last five position control points and three
orientation control points before the window start, plus the bias (30
coordinates).  A factor touching a leaving variable whose support
would reach past this boundary, a sliver at the mismatch of the two
grids, is dropped rather than absorbed, so the prior stays a marginal
over exactly that boundary.

\textbf{Re-centered prior.}
Before each solve, the prior is centered at the current warm start
$\mathbf{x}_0$, retaining its information matrix and setting its
gradient there to zero.  This is a curvature-only approximation to
the marginalized problem.  It generally changes the prior's mean,
so it does not preserve the full Gaussian marginal.  The prior still
penalizes subsequent departures from the warm start during the solve.

\textbf{Staleness of the prior.}
To reduce reliance on older linearizations during rotation, the
prior's position and orientation uncertainties are increased with
accumulated angular travel.  The update adds the variance $(k\theta)^2$ to every boundary position
and orientation coordinate of the marginal's covariance (a Woodbury
update, so the cross-covariances follow), with $\theta$ the rotation
accumulated over the marginalized stride and the rates $k = 0.12$\,m
and $0.01$\,rad per radian measured as the live-edge position and
orientation error increments per swept radian against ground truth
(the position rate is $0.12$, $0.12$ and $0.17$\,m/rad on the three
flights).  Angular travel is a proxy for linearization error, not a
direct measure of it.
Apart from this inflation, no global scale factor is applied to the
prior.  \autoref{sec:ablations} evaluates the effect of position
inflation, and \autoref{sec:nees} tests the output covariance.

\textbf{Live-edge warm start.}
The incoming trajectory segment is aligned to the previous solution
at its last data-supported position and orientation control points.  This
reduces mismatch at the transition before optimization.
